% ============================================================
%  计算机软件著作权登记申请表（正文）
% ============================================================
% 用法：直接修改下面各字段的第二个参数即可。勾选项用
%   \tjcrchk{...}    已勾选
%   \tjcrunchk{...}  未勾选
% 需要打印后手写的空白处用 \tjcrblank{宽度}。
% 字数限制等提示文字用 \tjcrred{...} 标红。
% “软件的主要功能”字数较多，已按句号/分号拆成若干续行
% （\tjcrcont），以便 longtable 在分页处自然断开；增删内容时
% 保持每条续行不超过本页高度即可。
% ============================================================

\vspace*{0.2cm}
{\centering\linespread{1}\fontsize{12bp}{18bp}\selectfont\heiti\bfseries
  计算机软件著作权登记申请表\par}
\vspace{0.5em}

\begin{tjcrform}

% ------------------------------------------------------------
% 权利取得方式
% ------------------------------------------------------------
\tjcrrow{权利取得方式}{%
  \linespread{1.5}\selectfont
  \tjcrchk{原始取得}\newline
  \tjcrunchk{继受取得}\quad(\tjcrunchk{受让}\quad\tjcrunchk{承受}\quad\tjcrunchk{继承})\newline
  \hspace*{2em}\tjcrunchk{软件是否已登记}\quad（原登记号：\tjcrblank{3cm}）\newline
  \hspace*{2em}\tjcrunchk{原登记是否做过变更或补充}\newline
  \hspace*{4em}（变更或补充证明编号：\tjcrblank{2.8cm}）}

% ------------------------------------------------------------
% 软件基本信息
% ------------------------------------------------------------
\tjcrsection{软件基本信息}

\tjcrrow{软件名称}{码极客具身智能体能力管理软件}
\tjcrrow{软件简称}{\tjcrgray{没有可以不填写}}
\tjcrrow{版本号}{V1.0}

\tjcrrow{权利范围}{%
  \linespread{1.5}\selectfont
  \tjcrchk{全部}\newline
  \tjcrunchk{部分}\newline
  \hspace*{2em}\tjcrunchk{发表权}\quad\tjcrunchk{署名权}\quad\tjcrunchk{修改权}\quad
  \tjcrunchk{复制权}\quad\tjcrunchk{发行权}\newline
  \hspace*{2em}\tjcrunchk{出租权}\quad\tjcrunchk{信息网络传播权}\quad
  \tjcrunchk{翻译权}\quad\tjcrunchk{应当由著作权人享有的其他权利}}

% ------------------------------------------------------------
% 软件开发信息
% ------------------------------------------------------------
\tjcrsection{软件开发信息}

\tjcrrow{软件分类}{%
  \tjcrchk{应用软件}\quad\tjcrunchk{嵌入式软件}\quad
  \tjcrunchk{中间件}\quad\tjcrunchk{操作系统}}

\tjcrrow{软件作品说明}{%
  \linespread{1.5}\selectfont
  \tjcrchk{原创}\qquad\tjcrunchk{修改（含翻译软件、合成软件）}\newline
  \hspace*{4em}修改合成或翻译说明（50字内）：\newline
  \tjcrunchk{该软件已登记}\quad\tjcrblank{3.2cm}\quad（上传原证书扫描件）\newline
  是修改、翻译或合成别人软件且需要原权利人授权的（上传授权书）}

\tjcrrow{开发方式}{%
  \linespread{1.5}\selectfont
  \tjcrchk{独立开发}\quad\tjcrunchk{合作开发}\quad\tjcrunchk{委托开发}\quad
  \tjunchecked\hspace{0.25em}下达任务开发（除独立开发，其他开发方式需要提供合同或协议书）}

\tjcrrow{开发完成日期}{2026.09.05}
\tjcrrow{发布状态}{\tjcrunchk{已发表}\qquad\tjcrchk{未发表}}

% “已发表信息”两小行放在同一单元格内，行间不画横线，标签自然居中。
\tjcrrow{已发表信息}{%
  \linespread{1.5}\selectfont
  首次发表日期：\tjcrblank{3cm}\newline
  首次发表地点：国家：中国\quad 城市：上海}

% ------------------------------------------------------------
% 著作权人（子表：表头 + 两行空白填写行）
% ------------------------------------------------------------
\tjcrrowraw{\tjcrstrut 著作权人\newline（业务部门\newline 只需填写\newline 姓名或名称）}{%
  \begin{tjcrcopyrighters}
    姓名或名称 & 类别 & 证件类型 & 证件号码 & 国籍 & 省份/\newline 城市 & 园区 \\ \hline
    & & & & & & \\ \hline
    & & & & & & \\ \hline
  \end{tjcrcopyrighters}%
  % 略增本行高度，使“软件功能与特点”分节行另起一页，避免孤行。
  \vspace*{0.5\baselineskip}}

% ------------------------------------------------------------
% 软件功能与特点
% ------------------------------------------------------------
\tjcrsection{软件功能与特点}

\tjcrrow{开发的硬件环境（\tjcrred{50字符内}）}%
  {x86-64处理器、8GB以上内存的台式机或笔记本}
\tjcrrow{运行的硬件环境（\tjcrred{50字符内}）}%
  {x86\_64架PC或笔记本电脑，Intel Core i5及以上处理器，8GB及以上内存}
\tjcrrow{开发该软件的操作系统（\tjcrred{50字符内}）}{Windows 11}
\tjcrrow{软件开发环境/开发工具（\tjcrred{50字符内}）}%
  {VS Code、Python、FastAPI、Chrome开发者工具}
\tjcrrow{该软件的运行平台/操作系统（\tjcrred{50字符内}）}%
  {Windows、Linux或macOS，配合Chrome或Edge浏览器}
\tjcrrow{软件运行支撑环境/支持软件（\tjcrred{50字符内}）}%
  {Chrome或Edge浏览器，Python3.12、FastAPI、Docker}

\tjcrrow{编程语言}{%
  \linespread{1.5}\selectfont
  \tjcrunchk{Assembly language}\quad\tjcrunchk{C}\quad\tjcrunchk{C\#}\quad
  \tjcrunchk{C++}\quad\tjcrunchk{Delphi/Object Pascal}\quad\tjcrunchk{Go}\quad
  \tjcrchk{HTML}\newline
  \tjcrunchk{Java}\quad\tjcrchk{JavaScript}\quad\tjcrunchk{MATLAB}\quad
  \tjcrunchk{Objective-C}\quad\tjcrunchk{PHP}\quad\tjcrunchk{PL/SQL}\quad
  \tjcrunchk{Perl}\quad\tjcrchk{Python}\quad\tjcrunchk{R}\quad
  \tjcrunchk{Ruby}\quad\tjcrunchk{SQL}\quad\tjcrunchk{Swift}\newline
  \tjcrunchk{Visual Basic}\quad\tjcrunchk{Visual Basic .Net}\newline
  如有需要，请输入其他编程语言：}

\tjcrrow{源程序量}{7131行}
\tjcrrow{开发目的（\tjcrred{50字符内}）}%
  {构建具身智能体Recipe能力统一管理与调试终端。}
\tjcrrow{面向领域/行业（\tjcrred{50字符内}）}%
  {具身智能、机器人巡检、智慧城市与工业运维。}

% ------------------------------------------------------------
% 软件的主要功能：按句号/分号拆成多条续行，便于跨页
% ------------------------------------------------------------
\tjcrrow{软件的主要功能\newline（\tjcrred{500-1300字内}）}{%
  Geek Mind具身智能体能力管理软件面向具身智能体、巡检机器人和工业运维场景，
  用于对Agent框架可调用的Recipe能力层进行集中展示、登记、管理、监控和调试。}
\tjcrcont{%
  软件提供工业运维后台与大屏驾驶舱结合的Web界面，可在未连接真实机器人、
  Zenoh总线或agent-server服务时进入离线演示模式，加载内置能力目录和模拟本体状态。}
\tjcrcont{%
  系统支持Recipe能力总览，统计注册能力、已接入能力、端侧能力、云端AI能力、
  规划中能力和异常检测能力；支持能力目录管理，用户可新增、编辑、删除、搜索和筛选能力，
  并维护能力名称、分类、状态、部署位置、调用协议、接口端点、输入输出Schema、
  依赖服务和标签等信息。}
\tjcrcont{%
  软件内置巡检技能、导航技能、操作技能、视觉分析、异常检测、语音交互、空间交互、
  云端AI和扩展能力等分类，预置本体回零、目标点导航、机械臂末端位姿控制、
  机械臂关节控制、视觉反应式抓取、示教动作执行、底层速度指令、安全预检查、
  急停中止和故障复位等机器人能力。}
\tjcrcont{%
  同时扩展烟火检测、烟雾检测、违规摆摊识别、群众聚集检测、打架斗殴识别、
  非法停车检测、垃圾堆放检测、安全帽识别、人员跌倒检测、危险区域闯入检测、
  积水检测、云端VLM场景理解、OCR文字识别、ASR语音识别、TTS语音合成、
  异常事件复核和巡检报告生成等AI能力接口。}
\tjcrcont{%
  在线模式下，软件可连接现有om1-agent-server，读取/health和/state接口，
  动态同步机器人可用verb、本体状态、后端类型、FAULT故障锁存和示教动作；
  接口调试中心可生成请求JSON，调用/verb、/dispatch、/abort、/reset等接口
  并展示CompletionVerdict回执。}
\tjcrlastcont{%
  软件通过能力状态区分已接入、降级、未连接、规划中、故障和停用，
  既可管理真实端侧小脑能力，也可登记尚未实现的云端单点AI能力，
  为具身智能体能力接入、运行监测和扩展治理提供统一终端。}

% ------------------------------------------------------------
% 软件的技术特点
% ------------------------------------------------------------
\tjcrrow{软件的技术特点\newline（\tjcrred{最多3个，标签可选}）}{%
  \linespread{1.5}\selectfont
  \tjcrunchk{APP}\quad\tjcrunchk{游戏软件}\quad\tjcrunchk{教育软件}\quad
  \tjcrunchk{金融软件}\quad\tjcrunchk{医疗软件}\quad\tjcrunchk{地理信息软件}\newline
  \tjcrunchk{云计算软件}\quad\tjcrunchk{信息安全软件}\quad
  \tjcrunchk{大数据软件}\quad\tjcrchk{人工智能软件}\quad
  \tjcrunchk{VR软件}\quad\tjcrunchk{5G软件}\quad\tjcrunchk{小程序}\quad
  \tjcrchk{物联网软件}\quad\tjcrchk{智慧城市软件}\par
  \vspace{0.4em}
  技术特点（100字内）：（必填）采用Recipe能力目录与前后端分层设计，
  支持离线演示、在线探测、能力增删改查、接口代理调试及机器人安全控制。}

\end{tjcrform}
